\subsubsection{Meaning of fixed routing}

The software does not store one canonical physical route per OD pair. It builds one time-expanded graph for the
complete timetable, with event nodes tied to scheduled trips and one origin centroid for every stop and departure-time
bin. Feasible alternatives are implicit link sequences in this graph rather than an enumerated path table. Access links
from a time-bin centroid reach only departures in its configured time window and carry the corresponding schedule-
deviation costs. Thus, the same physical OD in two time bins can face different lines and scheduled trips; even if its
physical stop or line sequence is unchanged, its trip sequence can differ.

For each destination, the Dial recursion produces one conditional probability per enabled time-expanded link. A demand
time bin is an aggregation interval: all demand in one OD-time cell is injected at the same time-bin centroid and shares
its distribution over accessible departure events. The current model contains no finer within-bin departure-cohort
index.

In the ordinary assignment call these destination value functions and link probabilities are recomputed before demand
is loaded. In \emph{fixed routing}, the effective link masks and probabilities are computed once for a fixed graph,
timetable, generalized-cost vector, destination masks, and dispersion parameter $\theta$, then reused for every demand
evaluation. Complete, sharded, concurrent, matrix-free, and gravity-model operators preserve this same routing contract.
The gravity model changes how OD-time demand is generated; it does not define different routing probabilities.

Consequently, ``the route set is constant throughout the day'' and ``route-choice probabilities are constant throughout
the day'' are false in general. ``Routing probabilities are held fixed during fixed-routing estimation'' is true.
Furthermore, routing is currently independent of demand and crowding: nominal ride-link capacities are stored as
metadata but do not cap boarding, generate denied boarding, change passenger propagation, or enter utility. Changing
only OD demand therefore leaves probabilities unchanged, and fixed loading agrees with a freshly recomputed assignment
when all routing-sensitive inputs are unchanged.

If crowding or capacity becomes endogenous, fixed routing would no longer be exact. A possible outer fixed-point scheme
would alternate demand estimation, capacity-aware dynamic assignment, routing-cache reconstruction, and re-estimation
until demand and assigned flows stabilize.


